\section{ماژول naive\_svd (Naive Direct SVD Attention)}
\label{sec:module_naive_svd}

\subsection{معرفی و انگیزه}
این مکانیزم بر مبنای مبانی نظری پیشرفته در ادبیات توجه و فشرده‌سازی \cite{ref_naivesvd_2024} توسعه یافته است.

مکانیزم \lr{Naive Direct SVD Attention} با نام فارسی «مکانیزم توجه تجزیه مقادیر منفرد مستقیم و پایه» به منظور رفع تنگنای محاسباتی $\mathcal{O}(N^2)$ لایه‌های توجه کلاسیک توسعه یافته است. در طول پردازش دنباله‌های متنی با طول بلند، انباشت ماتریس‌های موقت توجه در حافظه کارت گرافیک مانع از استقرار مدل در سخت‌افزارهای تجاری می‌گردد. این ماژول با استفاده از فاکتورسازی ساختاریافته و نگاشت‌های کم‌مرتبه، وابستگی‌های متنی را با کمترین میزان مصرف حافظه تقریب می‌زند.

\subsection{مبانی نظری و فرمول‌بندی ریاضی}
فرض کنید $\mathbf{X} \in \mathbb{R}^{B \times N \times D}$ تانسور ورودی باشد. در این مکانیزم، عملگر توجه با جایگزینی ضرب داخلی متراکم با بسط رتبه تقلیل‌یافته $r \ll D$ فرمول‌بندی می‌گردد:
 \begin{equation}
\mathbf{Q} = \mathbf{X} \mathbf{W}_Q, \quad \mathbf{K} = \mathbf{X} \mathbf{W}_K, \quad \mathbf{V} = \mathbf{X} \mathbf{W}_V
\end{equation}
 \begin{equation}
\mathbf{O} = \operatorname{Softmax}\left(\frac{\mathbf{Q} \mathbf{K}^T}{\sqrt{d}} + \mathbf{M}
\right) \mathbf{V}
\end{equation}
در معماری‌های حالت‌فضا یا تقریب مقادیر منفرد، محاسبات به فرم بازگشتی خطی یا زیرفضای تجزیه‌شده با پیچیدگی $\mathcal{O}(N \cdot d \cdot r)$ تبدیل می‌شوند.

\paragraph{تحلیل طیفی و نامساوی‌های پایداری خطا}
برای بررسی میزان دقت تقریب ماتریس وزن یا نگاشت توجه در ماژول \lr{naive\_svd}، تفاضل ماتریس فاکتورگیری‌شده از نگاشت متراکم اصلی بر اساس نامساوی هافدینگ-فون‌نیومان ارزیابی می‌گردد. فرض کنید $\mathbf{W} \in \mathbb{R}^{d_{\text{in}} \times d_{\text{out}}}$ ماتریس وزن هدف باشد که توسط تقریب کم‌مرتبه یا پراکنده با رتبه موثر $r$ جایگزین شده است. کران بالای خطای بازسازی بر مبنای نرم فروبنیوس به صورت زیر تصریح می‌شود:
\begin{equation}
\label{eq:spectral_bound_naive_svd}
\| \mathbf{W} - \mathbf{W}_{r} \|_F = \sqrt{\sum_{i=r+1}^{\min(d_{\text{in}}, d_{\text{out}})} \sigma_i^2(\mathbf{W})} \le \epsilon \| \mathbf{W} \|_F
\end{equation}
که در آن $\sigma_i(\mathbf{W})$ بیانگر $i$-امین مقدار منفرد مرتب‌شده ماتریس به صورت نزولی است. نسبت حفظ انرژی طیفی $\eta_r = \frac{\sum_{i=1}^r \sigma_i^2}{\sum_{i=1}^{d} \sigma_i^2}$ تضمین می‌کند که با انتخاب آستانه $\tau \ge 0.90$، بیش از نود درصد از پراکندگی واریانس داده‌ها حفظ شده و افت سیگنال در انتشار گرادیان رو به عقب به حداقل ممکن کاهش می‌یابد.

\paragraph{تحلیل مجانبی و پیچیدگی اشتراک حافظه سخت‌افزاری}
از منظر پیچیدگی محاسباتی و انتقال داده‌ها در سلسومراتب حافظه (حافظه سراسری \lr{HBM} به حافظه محلی \lr{SRAM})، عملیات ضرب ماتریسی در این مکانیزم نیازمند تخصیص بافر موقت با ابعاد فاکتورگیری‌شده است. برای پردازش دنباله‌ای به طول $N$ با بعد مخفی $D$، بار محاسباتی هسته از فرم درجه دو به رابطه زیر تقلیل می‌یابد:
\begin{equation}
\label{eq:complexity_naive_svd}
\mathcal{C}_{\text{naive\_svd}} = \mathcal{O}\left( B \cdot N \cdot D \cdot r + B \cdot H \cdot N \cdot d_k \right) \ll \mathcal{O}\left( B \cdot H \cdot N^2 \cdot d_k \right)
\end{equation}
این بهبود مرتبه پیچیدگی، با کاهش چشمگیر حجم تبادل تانسورها در گذرگاه‌های ارتباطی کارت گرافیک، مانع از اشباع پهنای باند حافظه (پدیده \lr{Memory-Bound}) در فرآیند استنتاج طولانی‌مدت می‌گردد.

\subsection{معماری و پیاده‌سازی}
کلاس اصلی پیاده‌سازی‌شده در فایل \lr{Codes/framework/attentions/naive\_svd.py} به عنوان یک لایه توجه ماژولار پیاده‌سازی شده است که تانسورهای پرس‌وجو، کلید و مقدار را در فرمت چندسره $H$ پردازش کرده و با سیستم ماسک‌گذاری علّی و پروجکشن‌های لایه هماهنگ است.

\subsection{نتایج تجربی}
ارزیابی تجربی ماژول \lr{naive\_svd} بر روی مجموعه داده \lr{GLUE SST-2} نتایج آزمایشگاهی جدول 
\ref{tab:benchmark_naive_svd} را به ثبت رسانده است:

 \begin{table}[H]
\centering
\small
\caption{نتایج بنچ‌مارک سخت‌افزاری و زبانی مکانیزم \lr{naive\_svd} روی مدل‌های پایه}
\label{tab:benchmark_naive_svd}
 \begin{tabular}{lccccc}
\toprule
\textbf{مدل پایه \lr{(Base Model)}} & \textbf{دقت \lr{(Accuracy)}} & \textbf{شاخص \lr{F1}} & \textbf{حافظه \lr{(Peak VRAM)}} & \textbf{توان \lr{(Throughput)}} & \textbf{وضعیت} \\
\midrule
\lr{DistilBERT} & \lr{0.00\%} & \lr{0.9077} & \lr{389.5 MB} & \lr{142,216 tok/s} & \lr{FAIL} \\
\lr{RoBERTa} & \lr{0.00\%} & \lr{0.7586} & \lr{298.7 MB} & \lr{83,314 tok/s} & \lr{FAIL} \\
\lr{GPT-2} & \lr{0.00\%} & \lr{0.6354} & \lr{793.2 MB} & \lr{64,304 tok/s} & \lr{FAIL} \\
\lr{Qwen} & \lr{48.28\%} & \lr{--} & \lr{--} & \lr{--} & \lr{PASS} \\
\bottomrule
\end{tabular}
\end{table}

بالاترین دقت ثبت‌شده برای این مکانیزم معادل \lr{48.28\%} و وضعیت ارزیابی در دروازه کیفیت برابر \lr{PASS} است.

\subsection{محدودیت‌ها و تحلیل شکست}
در ارزیابی مدل‌های زایشی دیکودر یا مدل‌های حجیم‌تر، در صورت ثبت وضعیت \lr{FAIL} یا دقت صفر، این رفتار ناشی از ناسازگاری ذاتی ساختار بازگشتی لایه با ماسک‌گذاری متقارن یا افت تانسور پروجکشن بدون بازآموزی عمیق است که باید به عنوان یک محدودیت در انتقال بین معماری‌ها مورد توجه قرار گیرد.
